\section{Experimental Setup}
\label{sec:experiments}

Our experiments test whether RLVR on procedurally generated \trace instances
improves visual reasoning beyond the distributions represented by the training
tasks. We initialize from Qwen2.5-VL-3B-Instruct and
Qwen2.5-VL-7B-Instruct~\citep{bai2025qwen25vl} and apply the same training
procedure at both model scales. Each resulting checkpoint is compared with its
corresponding base checkpoint under a common external-evaluation protocol.

\subsection{Training data and reward}

Both model scales are trained on the same fixed set of 64,000 instances,
comprising 64 independently generated examples from each of the 1,000 \trace
tasks. A separate held-out \trace validation set contains two previously unseen
instances per task, for 2,000 examples in total. All reported results use the
final checkpoint after 500 updates.

The model may produce an unrestricted reasoning trace, but its response must
terminate with a JSON object of the form
\texttt{\{"answer": ...\}}. Let $R_a \in \{0,1\}$ indicate whether the
submitted answer exactly matches the target after type-aware canonicalization,
and let $R_f \in \{0,1\}$ indicate whether the response ends with a valid JSON
object containing exactly the requested key. The reward is
\begin{equation}
  R = 0.95 R_a + 0.05 R_f.
\end{equation}
Answer correctness therefore determines the reward, while the smaller format
term encourages a consistently parseable output interface.

\subsection{Optimization}

We optimize both models with group relative policy optimization (GRPO). Each
update samples 128 prompts and generates eight responses per prompt. Rewards
are normalized within each eight-response group and optimized using a clipped
token-level policy objective. The 500 updates correspond to one shuffled pass
over the 64,000 training prompts and produce 512,000 sampled responses in
total.

We use full-parameter training with BF16 parameters, a learning rate of
$10^{-6}$, and one policy epoch per update. Rollout sampling uses temperature
1.0 and top-$p$ 1.0. Training is performed on eight H100 80GB GPUs and takes
12.2 hours for the 3B model and 13.9 hours for the 7B model. The two scales
share the same data, reward, sampling schedule, and optimization
hyperparameters; Appendix~\ref{app:experiments} provides the complete
configuration.

\subsection{External evaluation}

To measure transfer beyond the \trace task distributions, we evaluate each
base and trained checkpoint on 24 external benchmarks spanning six analysis
groups: charts and tables, visual mathematics, science and general reasoning,
spatial reasoning, perception and counting, and puzzles and logic. Each group
contains four benchmarks. Table~\ref{tab:main-results} reports the complete
per-benchmark results, while Appendix Table~\ref{tab:evaluation-suite} lists
the evaluated splits, sample counts, and benchmark-specific references.

We evaluate each designated split in full, yielding 32,805 examples per model
and decoding seed. Each benchmark retains its native evaluation metric.
Category scores are computed as unweighted means of the four benchmarks in
that group, and the overall score is the unweighted mean across all 24
benchmarks. This macro-averaging prevents benchmarks with larger evaluation
sets from dominating the aggregate result.

For each fixed checkpoint, we repeat decoding with three seeds and report the
mean and sample standard deviation. The seeds affect only response generation;
the model weights remain unchanged. All checkpoints are evaluated with the same
benchmark prompts, image-preprocessing bounds, parsers, scorers, and
VLMEvalKit version~\citep{duan2024vlmevalkit}. Exact inference settings are provided in
Appendix~\ref{app:experiments}.

At the 7B scale, we additionally evaluate publicly available synthetic-data
RLVR checkpoints from Game-RL, Sphinx, and
PC-GRPO~\citep{tong2025gamerl,alam2026sphinx,jeddi2025pcgrpo}. We report
Vero~\citep{sarch2026vero} separately as a real-image reference because its
600K-example training mixture draws from 59 datasets and includes real-image
data, whereas our runs use 64K procedurally generated instances. All
checkpoints are evaluated under the same external protocol. Because their
training data, optimization, and compute are not matched, these comparisons
characterize checkpoint outcomes rather than isolate the effect of a specific
training method or data source.
